% Discusión del corte de resultados documentado en ch5_results.tex.
% Referencias metodológicas: ch4_methodology_refactorizada.tex, activo en main.tex.
% Los cálculos adicionales son derivaciones de las tablas, del resumen verificado
% y de la implementación local del generador; no representan nuevas ejecuciones.
% Se conservan u_T y u_L donde la exportación no acredita unidades SI.

\chapter{Discusión de resultados}
\label{ch:discusion}

Los resultados del Capítulo~\ref{ch:resultados} permiten evaluar la coherencia y el grado de desarrollo de una cadena de observación, simulación y predicción orientada al control termo--higrométrico. Su contribución más sólida es la combinación de una preparación temporal auditable, pruebas conservativas de modelos físicos y comparaciones predictivas con una referencia explícita. La evidencia también delimita dónde termina esa demostración: el conjunto analizado contiene series sintéticas y pruebas numéricas, mientras que la validación independiente del recinto y la comparación en lazo cerrado permanecen pendientes. Esta distinción determina la interpretación científica de los errores, los balances y las posibilidades de aplicación.

La discusión relaciona cada hallazgo con el procedimiento que lo sustenta en el Capítulo~\ref{ch:metodologia} y con sus fundamentos en el Capítulo~\ref{ch:marco_teorico}. Además de interpretar las cifras reportadas, se desarrollan consecuencias analíticas sobre soporte espacial, ruido del generador, identificabilidad y capacidad de actuación. Estas derivaciones se distinguen de las ejecuciones originales y conservan las unidades y condiciones de validez de sus fuentes.

\section{Alcance de la evidencia y correspondencia entre modelos}
\label{dis:alcance}

\subsection{Qué se ha demostrado en cada nivel de la investigación}

La secuencia de la Tabla~\ref{met:tabla_fases} conecta caracterización, observación, verificación, identificación, predicción y control. Los resultados disponibles cubren esas fases con distinta profundidad. La integridad de la tabla de sensores acredita consistencia de los datos procesados; la cavidad examina una implementación de convección; los balances de la red cuantifican su cierre numérico; y el bloque cronológico final permite comparar predictores sobre una serie generada. Cada comprobación responde a una pregunta diferente. Un resultado favorable en una de ellas no sustituye los requisitos de las demás.

La Tabla~\ref{tab:dis_evidencia} organiza las conclusiones admisibles. Su lectura conjunta muestra que la investigación dispone de componentes computacionales evaluados, aunque todavía carece del enlace experimental que permitiría atribuir sus prestaciones a Coffee Solutions. En particular, la coincidencia con una referencia producida por un modelo verifica capacidad de reproducción dentro de ese marco; la validación física requiere observaciones independientes y un operador de medición compatible, conforme a las Secciones~\ref{sec:observacion_validez} y~\ref{met:validacion}.

\begin{table}[htbp]
\centering
\small
\caption{Nivel de evidencia alcanzado y alcance de las inferencias principales.}
\label{tab:dis_evidencia}
\begingroup
\renewcommand{\arraystretch}{1.25}
\setlength{\tabcolsep}{4pt}
\begin{tabularx}{\textwidth}{>{\raggedright\arraybackslash}p{0.20\textwidth}>{\raggedright\arraybackslash}X>{\raggedright\arraybackslash}X}
\toprule
Componente & Evidencia disponible & Inferencia respaldada \\
\midrule
Sensado y datos & Serie sintética completa y auditoría de identidades, tiempos y antecedentes. & Consistencia del conjunto procesado y corrección de su preparación causal. \\
CFD tridimensional & Arranque aire--vapor de $0{,}05$ s en recinto vacío. & Funcionamiento inicial de la ruta de solución y comportamiento psicrométrico compatible. \\
Cavidad & Dos Rayleigh, tres mallas y variación del paso; caso radiativo adicional. & Respuesta convectiva y sensibilidad numérica; incertidumbre espacial aún sin resolver. \\
Red reducida & Prueba de seis estados, balances y contraste de integración. & Coherencia conservativa y estabilidad de la solución del caso suministrado. \\
Árboles y ensambles & Prueba cronológica y comparación emparejada sobre datos generados. & Habilidad predictiva dentro del soporte del generador. \\
PINN & Aproximación de referencias numéricas de seis estados. & Fidelidad al problema reducido ensayado y residual físico evaluado. \\
Enriquecimiento y control & Procedimiento formulado; comparaciones comunes sin ejecutar. & Hipótesis de integración y evaluación definidas, con efecto aún indeterminado. \\
\bottomrule
\end{tabularx}
\endgroup
\end{table}

\subsection{Recinto vacío y representación con producto}

Existe una diferencia de alcance que debe hacerse explícita. La metodología integrada define el recinto sin carga y el estado $[T_c,T_b,T_a,W]^{\mathsf T}$ en la Ecuación~\eqref{met:estados_entradas}. En cambio, la red y la PINN exportadas en las Tablas~\ref{tab:res_red_estados} y~\ref{tab:res_pinn} utilizan $[T_c,T_b,T_p,T_a,W,M]^{\mathsf T}$, correspondiente a la representación con producto de la Sección~\ref{subsec:estado_termohigrometrico}. Los dos estados adicionales incorporan almacenamiento térmico y agua del material, con intercambios que modifican también la evolución del aire.

Por tanto, las pruebas de seis estados se interpretan como evidencia de una formulación ampliada y de su implementación. Sus residuales y errores no se trasladan automáticamente a la red de cuatro estados de la metodología vigente. Establecer esa correspondencia requiere una configuración sin producto, retirar sus intercambios y fuentes, y repetir balances y contrastes sobre la formulación resultante. Eliminar únicamente $T_p$ y $M$ de una tabla conservaría trayectorias afectadas por una carga que la planta metodológica excluye.

Esta separación permite reconocer el valor del trabajo realizado sin confundir configuraciones. El recinto vacío constituye una base para identificar captación, pérdidas y renovación antes de introducir la complejidad del producto. La representación ampliada ofrece una vía para estudiar secado posteriormente, pero su evaluación actual no acredita temperatura, humedad ni calidad de un lote real de café.
\FloatBarrier

\section{Calidad de los datos, metrología y representatividad espacial}
\label{dis:observacion}

\subsection{Integridad temporal y funcionamiento de la adquisición}

Los $466\,560$ registros y la cobertura nominal del $100\%$ de la Tabla~\ref{tab:res_cobertura} permiten construir ventanas regulares sin introducir imputaciones por ausencia de canales. La recuperación de $51\,836$ antecedentes por canal, documentada mediante la Ecuación~\eqref{eq:res_antecedentes_posterior}, tiene una consecuencia espacial concreta: evita excluir sistemáticamente las cuatro posiciones posteriores al preparar historias predictivas. La depuración conserva así una superficie cuya representación ya es menor que la de las laterales.

Esta corrección ilustra por qué el contrato de datos de la Sección~\ref{met:arquitectura} forma parte del método científico. Un identificador nulo válido puede convertirse en una pérdida de información físicamente localizada si el agrupamiento lo interpreta como ausencia de identidad. El error resultante afectaría los pesos espaciales y las poblaciones de entrenamiento, aun cuando todos los valores de temperatura y humedad fueran numéricamente admisibles.

La regularidad de $20$ s caracteriza el archivo generado. Evaluar la plataforma ESP32--nube de la Figura~\ref{met:fig_arquitectura} exige relacionar conversiones, mensajes emitidos, recepciones y retransmisiones. La cobertura de una tabla final no identifica retrasos, lecturas retenidas o pérdidas recuperadas. Tampoco permite juzgar si el período nominal resuelve los transitorios del recinto: esa decisión requiere conocer respuesta instrumental, constantes de tiempo y retardos de actuación, como establece la Sección~\ref{met:campana}.

\subsection{Concordancia térmica y dependencia entre canales}

La diferencia media DS--DHT de $1{,}4999\,u_T$ y su desviación próxima a $0{,}100\,u_T$ reproducen el desplazamiento y el ruido configurados en el generador (Tabla~\ref{tab:res_concordancia}). La estructura de generación explica esta regularidad mediante
\begin{equation}
T_{\mathrm{DHT}}=T_{\mathrm{DS}}-1{,}5+\varepsilon_D,
\qquad
T_{\mathrm{DS}}-T_{\mathrm{DHT}}=1{,}5-\varepsilon_D.
\label{eq:dis_concordancia_generador}
\end{equation}
La variación común de ambos canales se cancela en la resta. En consecuencia, una dispersión pequeña de la diferencia puede coexistir con errores comunes importantes respecto de una temperatura física desconocida. La concordancia caracteriza la construcción relativa de las señales; establecer exactitud requiere la referencia independiente y el presupuesto de incertidumbre definidos en la Sección~\ref{met:sensado}.

Tampoco procede aplicar $1{,}5$ como corrección al DHT11 instalado. La exportación conserva una escala térmica $u_T$ sin confirmación metrológica, y el desplazamiento fue impuesto por software. Del mismo modo, un RMSE predictivo inferior a la exactitud nominal de un instrumento describiría aproximación de sus lecturas, sin demostrar una temperatura real más exacta. Resolución, incertidumbre instrumental y error de pronóstico intervienen en problemas distintos.

La dependencia entre canales también afecta la interpretación multivariable. El DHT térmico hereda la realización DS, y la humedad generada utiliza a su vez el canal DHT. Las tres salidas contienen información compartida por construcción. Una mejora simultánea constituye evidencia de desempeño sobre esa tarea conjunta, aunque no equivale a tres verificaciones físicas independientes del microclima.

\subsection{Heterogeneidad observada y significado del promedio}

La diferencia térmica media derecha--izquierda de aproximadamente $2{,}000\,u_T$ y las desviaciones espaciales intrapared próximas a $0{,}32\,u_T$ muestran que una dispersión pequeña dentro de cada pared puede coexistir con diferencias relevantes entre paredes (Tablas~\ref{tab:res_descriptivos} y~\ref{tab:res_uniformidad}). Comparar ambas escalas no establece una razón exacta de varianzas, porque una corresponde a medias temporales y otra a dispersión simultánea. Sí revela que evaluar únicamente la uniformidad intrapared omite parte de la estructura espacial.

La relación entre esas componentes se expresa, para cada instante y pesos coherentes, mediante
\begin{equation}
\sigma_{q,\mathrm{total}}^2(t)=
\sum_s\pi_s\sigma_{q,s}^2(t)
+\sum_s\pi_s\bigl[\bar q_s(t)-\bar q(t)\bigr]^2,
\qquad \sum_s\pi_s=1.
\label{eq:dis_varianza_espacial}
\end{equation}
El primer término representa heterogeneidad interna de las superficies y el segundo, separación entre sus promedios. Esta descomposición desarrolla el indicador de la Ecuación~\eqref{met:dispersion_espacial} y muestra qué información se pierde al utilizar un único promedio. Las tablas agregadas disponibles no permiten reconstruir toda su evolución temporal.

La ponderación por registros asigna $16/36$ a cada lateral y $4/36$ a la pared posterior. Esos pesos responden al diseño de observación. Incluso una ponderación por las áreas nominales de la Tabla~\ref{tab:res_geometria} describiría las tres superficies representadas, sin incorporar techo, base, pared frontal ni distribución volumétrica del aire. Por ello, el agregado de $36{,}161\,u_T$ no identifica directamente $T_a$ de la red ni una temperatura media de exposición del café. La comparación requiere el operador de la Sección~\ref{met:operadores}.

La excedencia geométrica máxima de $1{,}106\times10^{-3}\,u_L$ parece pequeña frente a las dimensiones nominales, pero su importancia depende del gradiente y de la frontera donde se interpole. Una ubicación apenas exterior puede producir un valor indefinido o tomar una región sólida. La transformación de la Ecuación~\eqref{met:coordenadas_paredes} proporciona el procedimiento de enlace; su aplicación a esta exportación queda condicionada a acreditar unidades y montaje. Corregir coordenadas mediante recorte automático podría ocultar una inconsistencia sistemática.

Finalmente, la mayor temperatura derecha y la menor humedad relativa de esa pared pertenecen al patrón impuesto por el generador. La Figura~\ref{fig:res_horarias} no permite atribuir esas diferencias a orientación solar, aberturas o ventiladores reales. Tales mecanismos son hipótesis contrastables mediante meteorología y acciones sincronizadas, cuya ausencia limita la atribución causal.
\FloatBarrier

\section{Interpretación termo--higrométrica y alcance respecto del secado}
\label{dis:psicrometria}

\subsection{Humedad relativa, calentamiento y extracción de agua}

La distinción entre humedad relativa y razón de humedad es decisiva para interpretar el caso CFD. La Tabla~\ref{tab:res_cfd} muestra un aumento de temperatura media de $0{,}012322$ K y una disminución de humedad relativa de $0{,}015719$ puntos porcentuales, con $Y$ y $W$ constantes. A presión constante, la relación de la Ecuación~\eqref{met:conversion_humedad} conduce a
\begin{equation}
p_v=\frac{pW}{\varepsilon_m+W},\qquad
\left.\frac{\partial\ln\phi}{\partial T}\right|_{p,W}
=-\frac{\mathrm d\ln p_{ws}}{\mathrm dT}<0.
\label{eq:dis_humedad_calentamiento}
\end{equation}
El calentamiento eleva la presión de saturación y reduce $\phi$ aunque el contenido de vapor por masa de aire seco permanezca invariable. El transitorio es compatible con ese mecanismo. La tabla no permite identificar todos los flujos de frontera, por lo que la constancia de composición tampoco demuestra por sí sola que el recinto esté sellado.

Esta interpretación impide utilizar una disminución de humedad relativa como medida de agua removida. Un regulador podría reducir $\phi$ mediante calentamiento y cumplir una referencia aparente sin evacuar vapor. De ahí la conveniencia de regular temperatura y $W$, según la Sección~\ref{met:control}, y conservar $\phi$ como indicador termodinámico complementario. La presión y la compatibilidad térmica ausentes en la serie de sensores impiden comprobar experimentalmente esa transformación en el corte actual.

\subsection{Capacidad de la ventilación y efecto de introducir producto}

El balance de humedad del recinto vacío, Ecuación~\eqref{met:red_humedad}, permite deducir una restricción operativa. Para inventario constante, flujo exterior $F_{\mathrm{ext}}>0$ y ambiente constante,
\begin{equation}
W(t)=W_{\mathrm{amb}}+
\bigl[W(0)-W_{\mathrm{amb}}\bigr]
\exp\left(-\frac{F_{\mathrm{ext}}}{m_{da}}t\right).
\label{eq:dis_relajacion_humedad}
\end{equation}
La ventilación acelera la aproximación a $W_{\mathrm{amb}}$ y no puede sostener, mediante ese mecanismo aislado, un valor inferior al exterior. Una referencia higrométrica debe respetar esta limitación. La humedad relativa exterior, considerada sin su temperatura y presión, no determina el sentido de intercambio de vapor.

La recirculación puede modificar mezcla y conductancias térmicas, mientras que la renovación exterior introduce el término de intercambio másico. Esta separación, establecida en la Sección~\ref{met:caudal}, resulta especialmente relevante para los equipos de la Tabla~\ref{tab:res_actuadores}: potencia, diámetro o velocidad nominal no identifican caudal efectivo ni fracción de renovación. Sin anemometría o un cierre presión--caudal contrastado, no puede establecerse cuál modo extrae más vapor ni cuál mejora la uniformidad.

Al introducir café aparece una fuente de evaporación. Bajo las mismas aproximaciones de inventario, el balance de la Ecuación~\eqref{eq:balance_humedad_simplificado} sería $m_{da}\dot W=F_{\mathrm{ext}}(W_{\mathrm{amb}}-W)+\dot m_e$. Un régimen aproximadamente estacionario exigiría
\begin{equation}
W-W_{\mathrm{amb}}\simeq\frac{\dot m_e}{F_{\mathrm{ext}}}.
\label{eq:dis_carga_ventilacion}
\end{equation}
Esta expresión condicional vincula carga evaporativa y capacidad de renovación: mantener una misma banda con mayor evaporación requiere mayor intercambio exterior, siempre que las condiciones térmicas permitan sostenerlo. El producto también introduce consumo latente, resistencia al flujo y almacenamiento, por lo que el ajuste del recinto vacío necesita una ampliación física antes de representar secado.

El intervalo numérico de $M$, de $0{,}1$ a $0{,}09477986$ kg/kg de materia seca, equivale a una diferencia de $0{,}00522014$ kg/kg y al $5{,}22\%$ del valor inicial (Tabla~\ref{tab:res_red_estados}). No corresponde a $5{,}22$ puntos porcentuales de humedad en base húmeda. Sin conocer la masa seca tampoco permite calcular kilogramos de agua removida. Además, su carácter numérico impide atribuirlo a un lote observado. La conversión de bases de la Sección~\ref{subsec:humedad_producto} y la necesidad de seguir $T_p$ delimitan cualquier inferencia sobre cinética, tiempo de secado o calidad.
\FloatBarrier

\section{CFD tridimensional: funcionamiento inicial y fidelidad pendiente}
\label{dis:cfd}

El código de salida cero y la evolución de los campos de la Sección~\ref{res:cfd} muestran que el caso aire--vapor puede ejecutarse y producir salidas interpretables. La conservación de una composición uniforme durante un calentamiento pequeño constituye una comprobación inicial útil de la conexión entre campos y posprocesamiento psicrométrico. Sin embargo, un campo de vapor uniforme apenas exige al tratamiento de gradientes, difusión y transporte entre corrientes con composiciones diferentes. Su constancia no somete a prueba toda la física higrométrica prevista en las Ecuaciones~\eqref{met:cfd_vapor} y~\eqref{met:cfd_energia}.

Los seis instantes cubren $0{,}05$ s, frente a una adquisición nominal cada $20$ s y a procesos térmicos cuya escala debe identificarse. Esa diferencia impide relacionar directamente el arranque con un episodio operativo, reconocer régimen estacionario o establecer respuesta a la actuación. La proximidad de las temperaturas inicial y final informa sobre ese intervalo; no establece que los gradientes del recinto sean pequeños durante su operación diurna.

Las Figuras~\ref{met:fig_cfd_dominio} y~\ref{met:fig_cfd_malla} definen una organización geométrica y de mallado. Evaluar su fidelidad exige la malla efectivamente ejecutada, calidad, resolución en aberturas y paredes, historial de convergencia y refinamientos compatibles. El estudio bidimensional no proporciona la incertidumbre de discretización del caso tridimensional, porque cambian geometría, fronteras, malla y mecanismos dominantes. Tampoco un balance global pequeño demostraría por sí solo exactitud de chorros, máximos locales o zonas de bajo movimiento.

En consecuencia, el papel mejor sustentado para este CFD es el de plataforma inicial para generar evidencia espacial, una vez completados los controles de la Sección~\ref{met:refinamiento}. Su uso posterior podría ayudar a identificar regiones donde el promedio de la red omite heterogeneidad importante y a diseñar operadores para el mapa de sensores. En el corte disponible, esa utilidad es una capacidad metodológica por contrastar. No se han exportado campos y mediciones independientes suficientes para determinar estratificación real, eficiencia de mezcla o distribución del flujo de Coffee Solutions.
\FloatBarrier

\section{Cavidad bidimensional: transferencia de calor y error numérico}
\label{dis:cavidad}

\subsection{Consistencia convectiva y alcance del refinamiento}

En la malla de $65\times65$, el aumento de $Ra$ de $10^3$ a $10^4$ eleva el Nusselt medio de $1{,}114170$ a $2{,}202040$, casi al doble, y aumenta los extremos de velocidad (Tabla~\ref{tab:res_cavidad_espacial}). Dentro de la cavidad idealizada, esta respuesta es compatible con una mayor contribución de la circulación a la transferencia entre paredes. Ambos Nusselt superan la escala conductiva unitaria del problema, aunque ello no convierte sus valores en coeficientes directamente transferibles al secador abierto.

La coincidencia caliente--fría del orden de $10^{-15}$ aporta consistencia integral en el caso sin radiación. Su precisión próxima al redondeo debe interpretarse junto con la discretización y las simetrías del problema: un código puede equilibrar los flujos calculados y conservar un error apreciable en ambos. Asimismo, resolver exactamente el sistema discreto de Poisson mediante transformada seno elimina una contribución algebraica de ese subproblema, pero conserva el error de aproximar el continuo por una malla.

La Figura~\ref{fig:res_refinamiento} muestra que el Nusselt cambia de sentido al refinar. Las diferencias entre las dos mallas más finas, $0{,}146\%$ y $0{,}851\%$, son medidas de sensibilidad y no cotas de error. La ausencia de una secuencia monótona acreditada impide aplicar sin examen la forma simple de extrapolación de la Ecuación~\eqref{met:richardson}. El procedimiento de estimación de incertidumbre exige documentar el comportamiento de la secuencia y revisar casos oscilatorios o diferencias próximas a cero \cite{celik2008}.

El costo de pasar de $17\times17$ a $65\times65$ crece aproximadamente $32$ veces a $Ra=10^3$ y $47$ veces a $Ra=10^4$, a partir de los tiempos de la Tabla~\ref{tab:res_cavidad_espacial}. Como también cambian los pasos y los tiempos de detención, esas razones no identifican una ley de complejidad del algoritmo. Sí muestran que aumentar resolución tuvo un costo medible sin cerrar todavía la estimación de incertidumbre espacial. La decisión de malla debe apoyarse en un presupuesto de error para las magnitudes de interés.

\subsection{Separación entre error espacial, temporal y de detención}

Las diferencias temporales de Nusselt del orden de $10^{-7}$ o menor son mucho más pequeñas que las variaciones entre mallas. Bajo las condiciones ensayadas, este contraste sugiere priorizar el examen espacial y del criterio de detención antes de reducir nuevamente el paso. La comparación no constituye una descomposición completa del error: los tiempos finales difieren y la sensibilidad de una magnitud integrada puede ser menor que la de velocidades o gradientes locales.

Los órdenes aparentes de $0{,}2693$ y $2{,}9231$ tampoco caracterizan de forma robusta el integrador. Cuando las diferencias entre soluciones son muy pequeñas, su cociente puede ser dominado por detención, cancelación o redondeo. Para atribuir un orden temporal se requiere comparar tiempos comunes o estados estacionarios suficientemente acreditados, conforme a la Sección~\ref{met:refinamiento}, y comprobar que las demás contribuciones sean menores.

El residual estacionario inferior a $10^{-5}$ indica cumplimiento del criterio configurado. No proporciona un error de temperatura, velocidad o Nusselt respecto de la solución continua. De igual modo, la advección de primer orden ofrece una posible contribución al error espacial, pero las tablas no permiten atribuirle por sí sola la oscilación observada. El examen de fronteras, evaluación de flujos, tolerancias y solución de referencia debe preceder a cualquier explicación causal del patrón de refinamiento.

\subsection{Radiación y redistribución del intercambio}

La disminución del Nusselt caliente de $1{,}112544$ a $1{,}082202$, equivalente a $2{,}727\%$, demuestra que activar el intercambio radiativo modifica el resultado convectivo del caso considerado (Tabla~\ref{tab:res_cavidad_radiacion}). El signo no implica menor transferencia total ni menor eficiencia de un secador. Convección y radiación comparten condiciones térmicas y pueden redistribuir sus contribuciones, como se desarrolla en la Sección~\ref{sec:radiacion_solar}.

La diferencia caliente--fría de $-0{,}024357$ exige contabilizar todos los intercambios y el almacenamiento con una misma convención de signos:
\begin{equation}
\frac{\mathrm dE}{\mathrm dt}=
\sum_{\Gamma_i}\left(
\dot Q_{i,\mathrm{conv/cond}}+\dot Q_{i,\mathrm{rad}}\right).
\label{eq:dis_balance_radiativo}
\end{equation}
En régimen estacionario el miembro izquierdo se anula para el dominio definido. La diferencia entre dos contribuciones convectivas solo puede interpretarse después de incluir las demás superficies y mecanismos. La tabla disponible no permite resolver ese cierre completo.

Además, el único caso radiativo utiliza emisividad $0{,}1$, longitud $0{,}1$ m y diferencia térmica $10$ K. Su resultado no establece una regla general para una cubierta de secador. El aporte defendible es mostrar sensibilidad al acoplamiento y justificar que la captación solar de onda corta y el intercambio térmico de onda larga se caractericen por separado. Para ampliar la inferencia se necesita comprobar balances y explorar propiedades y fronteras pertinentes a la envolvente real.
\FloatBarrier

\section{Red reducida: conservación, simplificación e identificabilidad}
\label{dis:red}

\subsection{Jerarquía de errores y cierre de balances}

La red ofrece la evidencia conservativa más explícita del conjunto: los residuales normalizados máximos de aire seco, agua y energía son $3{,}858\times10^{-4}$, $8{,}314\times10^{-5}$ y $6{,}032\times10^{-6}$ (Tabla~\ref{tab:res_red_balances}). Estas cantidades corresponden, respectivamente, a $0{,}03858\%$, $0{,}008314\%$ y $0{,}0006032\%$ de los denominadores definidos para cada residual. Los porcentajes describen cierre respecto de esas escalas; no representan incertidumbre física ni una fracción universal del inventario inicial.

El residual energético absoluto de $0{,}462455$ J adquiere significado mediante la normalización de la Ecuación~\eqref{met:residual_normalizado}. Los máximos absoluto y normalizado pueden ocurrir en instantes diferentes, de modo que su cociente no reconstruye una potencia o una energía de referencia única. Un cierre pequeño respalda la implementación de los balances elegidos, aunque un modelo con cierres constitutivos inadecuados también puede conservar exactamente las cantidades que representa.

La Tabla~\ref{tab:res_red_comparacion} permite ordenar algunas contribuciones numéricas. Para $T_a$, la diferencia entre inventarios es $9{,}279\times10^{-4}$ K, frente a $5{,}686\times10^{-7}$ K entre Radau y BDF y $9{,}347\times10^{-10}$ K al refinar Radau. La diferencia de formulación supera aproximadamente $1{,}6\times10^3$ veces la discrepancia entre integradores y cerca de $10^6$ veces la del refinamiento. Estas razones son comparaciones del caso, sin convertir ningún contraste en error respecto de una solución exacta.

El resultado indica que reducir aún más las tolerancias tendría una utilidad limitada antes de examinar parámetros, fronteras y simplificaciones. La coincidencia entre integradores verifica estabilidad de la solución implementada, mientras que la cercanía entre inventarios sugiere una reducción potencialmente útil en ese régimen. Su admisibilidad operativa depende del presupuesto de error definido antes del contraste y de episodios con calentamiento y renovación distintos.

La mínima temperatura del producto, $298{,}076397$ K, es compatible con el término evaporativo de la Ecuación~\eqref{eq:nodo_producto}. Para identificar la contribución causal harían falta flujos y trayectorias sincronizadas o una comparación controlada del término. Los extremos agregados, por sí solos, no establecen un balance instantáneo ni confirman la magnitud real del enfriamiento evaporativo.

\subsection{Recuperación escalar e identificación de la planta}

La recuperación de $G_{ce}=5{,}0000$ W/K demuestra que el procedimiento resuelve el problema condicional ensayado. El rango uno y el número de condición uno describen un Jacobiano con un único parámetro libre. En esa situación, un número de condición unitario no informa sobre correlaciones con capacidades, captación, ventilación u otras conductancias, porque estas cantidades se mantuvieron fijadas.

La dificultad de ampliar el ajuste se aprecia en un balance simplificado,
\begin{equation}
C\dot T=\alpha_{\mathrm{sol}}AG(t)-G_e(T-T_{\mathrm{amb}}).
\label{eq:dis_identificabilidad}
\end{equation}
Con excitación y observación suficientes, la trayectoria puede informar sobre las razones $\alpha_{\mathrm{sol}}A/C$ y $G_e/C$. Si los tres factores son desconocidos, multiplicar $C$, $\alpha_{\mathrm{sol}}A$ y $G_e$ por una misma constante conserva la dinámica. Este ejemplo no demuestra una deficiencia estructural de toda la red, pero explica por qué un buen ajuste de temperatura puede coexistir con parámetros físicamente ambiguos.

La identificación de la Sección~\ref{met:identificacion} necesita, por tanto, información complementaria: capacidades y áreas caracterizadas, irradiancia y ambiente sincronizados, enfriamientos para pérdidas y cambios de modo para renovación. El Jacobiano escalado y los perfiles de parámetros deberían examinarse en esa configuración conjunta. Sin esos contrastes, no es posible interpretar las cifras ajustadas como propiedades verificadas de la planta ni propagar de manera fundada su incertidumbre hacia control.

La red concentra el estado y permite estudiar transitorios a menor complejidad espacial, pero esa reducción también elimina gradientes. Su relación con el CFD requiere preservar inventarios y energía mediante la Sección~\ref{met:operadores}. Un ajuste obtenido comparando directamente un nodo uniforme con una lectura local podría absorber en una conductancia el error de soporte espacial. Ese parámetro sería efectivo para la comparación utilizada y su transferencia a otra distribución de sensores o ventilación necesitaría comprobarse.
\FloatBarrier

\section{Predicción temporal: magnitud de la mejora y estructura del generador}
\label{dis:prediccion}

\subsection{Qué explica la diferencia frente a persistencia}

La Figura~\ref{fig:res_errores} presenta el contraste predictivo más claro del estudio. Al pasar de $20$ a $1200$ s, el RMSE DS de persistencia crece aproximadamente $77{,}1\%$, mientras que el de LightGBM aumenta cerca del $1{,}64\%$, calculado con las cifras no redondeadas del resumen verificado. A $1200$ s, las reducciones reportadas son $53{,}21\%$ en DS con LightGBM, $49{,}77\%$ en DHT térmico con XGBoost y $42{,}65\%$ en humedad con XGBoost (Tabla~\ref{tab:res_sesgo_mae}). La ganancia es consistente con la capacidad de los ensambles para aproximar un patrón temporal y espacial que persistencia mantiene fijo en el último valor observado.

La estabilidad del error de los ensambles tiene, sin embargo, una interpretación condicionada. La serie de tres días repite ciclos diarios y desplazamientos espaciales; el bloque posterior continúa bajo las reglas del mismo generador. El contraste temporal evita mezclar arbitrariamente pasado y futuro, pero no introduce por sí solo cambios de meteorología, carga, configuración o funcionamiento instrumental. El resultado respalda generalización temporal dentro de ese soporte. La robustez ante transiciones no representadas necesita otras trayectorias.

La respuesta del árbol refuerza la necesidad de una referencia simple. En DS empeora frente a persistencia a $20$ y $300$ s, y los intervalos exploratorios conservan signo negativo (Tabla~\ref{tab:res_intervalos}). La partición del espacio de variables y el ajuste por hojas pueden perder precisión frente a una señal de corto plazo casi persistente. Ese mecanismo es una explicación plausible, sin una ablación que permita atribuir el deterioro exclusivamente a profundidad, tamaño de hoja o selección de variables. La complejidad del predictor debe justificarse por tarea y horizonte.

\subsection{Escala de error esperable a partir del ruido sintético}
\label{dis:ruido}

La implementación del generador disponible en el proyecto permite profundizar la interpretación de los RMSE. El canal DS suma a un ciclo diario y a efectos de pared y altura un ruido normal de desviación $0{,}15\,u_T$; el DHT térmico añade un desplazamiento y ruido de $0{,}10\,u_T$; la humedad añade ruido de un punto porcentual y depende de DHT con factor $0{,}8$.
\footnote{Derivación a partir de las funciones del generador en \texttt{temperature\_model.py} y \texttt{dht\_model.py}, dentro de la tubería \texttt{synthetic\_data}, y de la configuración \texttt{synthetic\_data} en \texttt{conf/base/parameters.yml}. Se utiliza la implementación disponible para explicar la estructura de la serie; no se ha realizado una nueva simulación ni se atribuyen estas escalas a sensores físicos.}
Se trata de información estructural del caso sintético y no de un presupuesto metrológico.

Para una salida $q(t)=\mu_q(t)+\varepsilon_q(t)$, si la innovación futura es independiente de la información causal $\mathcal I_t$, tiene media cero y varianza $\sigma_q^2$, cualquier pronóstico causal satisface en población
\begin{equation}
\begin{split}
\mathbb E\bigl[(\widehat q_{t+h}-q_{t+h})^2\bigr]
&=\mathbb E\bigl[(\widehat q_{t+h}-\mu_q(t+h))^2\bigr]+\sigma_q^2\\
&\geq\sigma_q^2.
\end{split}
\label{eq:dis_suelo_ruido}
\end{equation}
La igualdad requiere predecir exactamente la componente determinista. Esta cota del error cuadrático esperado depende de las hipótesis indicadas y no es un límite determinista para cada RMSE muestral.

Bajo independencia de las innovaciones de generación, y en la región de humedad sin recorte en $0$ o $100\%$, las escalas correspondientes son
\begin{align}
\sigma_{\mathrm{DS}}&=0{,}15\,u_T,\nonumber\\
\sigma_{\mathrm{DHT}}&=\sqrt{0{,}15^2+0{,}10^2}\,u_T
\simeq0{,}1803\,u_T,\nonumber\\
\sigma_{\phi}&=\sqrt{1^2+0{,}8^2(0{,}15^2+0{,}10^2)}
\simeq1{,}0103\ \text{puntos porcentuales}.
\label{eq:dis_escalas_ruido}
\end{align}
Los factores se interpretan en las unidades numéricas del generador. La propagación refleja que el ruido DS participa también en DHT y que ambos contribuyen a la humedad generada. Los extremos observados de humedad, $30{,}610$ y $79{,}017\%$, se encuentran alejados de los límites de recorte.

\begin{table}[htbp]
\centering
\small
\caption{Escalas analíticas de ruido del generador y errores predictivos observados.}
\label{tab:dis_ruido}
\setlength{\tabcolsep}{5pt}
\begin{tabular}{lrrrr}
\toprule
Canal & $\sigma_q$ & $\sqrt{2}\sigma_q$ & Persistencia, $20$ s & Ensamble, $1200$ s \\
\midrule
$T_{\mathrm{DS}}$ & $0{,}1500$ & $0{,}2121$ & $0{,}2122$ & $0{,}1758$ \\
$T_{\mathrm{DHT}}$ & $0{,}1803$ & $0{,}2550$ & $0{,}2551$ & $0{,}2014$ \\
$\phi_{\mathrm{DHT}}$ & $1{,}0103$ & $1{,}4288$ & $1{,}4169$ & $1{,}0622$ \\
\bottomrule
\end{tabular}
\par\vspace{3pt}
\begin{minipage}{0.98\textwidth}\footnotesize
Temperaturas en $u_T$ y humedad en puntos porcentuales. Las dos primeras columnas son derivaciones condicionales del generador. Los RMSE proceden de la Tabla~\ref{tab:res_rmse}; se utilizan LightGBM para DS y XGBoost para los otros dos canales. $\sqrt{2}\sigma_q$ representa la contribución de dos innovaciones independientes a persistencia, sin incluir la evolución determinista.
\end{minipage}
\end{table}

La Tabla~\ref{tab:dis_ruido} aporta dos interpretaciones. Primero, para persistencia y una trayectoria determinista dada,
\begin{equation}
\mathbb E\bigl[(q(t)-q(t+h))^2\bigr]
=\bigl[\mu_q(t)-\mu_q(t+h)\bigr]^2+2\sigma_q^2.
\label{eq:dis_ruido_persistencia}
\end{equation}
A corto plazo, la contribución determinista puede ser pequeña; los errores térmicos de $0{,}2122$ y $0{,}2551\,u_T$ concuerdan estrechamente con $\sqrt{2}\sigma_q$. Una parte de la ventaja de un predictor que suaviza la señal proviene, por tanto, de evitar trasladar el ruido de la última lectura al pronóstico. Este comportamiento no requiere haber aprendido balances termodinámicos.

Segundo, a $1200$ s los RMSE de los ensambles seleccionados para esta comparación superan las escalas $\sigma_q$ aproximadamente $17{,}2\%$, $11{,}7\%$ y $5{,}14\%$ en DS, DHT y humedad. La cercanía ayuda a explicar por qué el error cambia poco con el horizonte dentro de un caso periódico sencillo. No establece que se haya alcanzado un óptimo estadístico ni que esas cotas describan la planta real. La irreducibilidad es condicional a innovaciones futuras independientes, preparación causal y coincidencia entre la implementación examinada y la fuente de los datos.

La comparación aconseja incorporar, en desarrollo, una referencia de ciclo diario con efectos de pared y altura. Un modelo armónico o una regresión temporal sencilla permitiría separar cuánto de la ganancia se debe a conocer la periodicidad y cuánto añade el ensamble. También sería informativo retirar sucesivamente hora, posición y canales relacionados, conservando las particiones. Estas ablaciones no están ejecutadas en el conjunto exportado y no se les asigna una ganancia anticipada.

\subsection{Sesgo, horizontes y generalización espacial}

Los sesgos reducidos de los ensambles a $1200$ s indican que el error medio global es pequeño en el bloque final. El RMSE sigue siendo apreciable porque promedia error cuadrático, y un sesgo cercano a cero puede resultar de cancelaciones entre fases del día o posiciones. La ausencia de sesgo global no garantiza límites de error en zonas calientes, transiciones o extremos de humedad. Los cuantiles y desgloses por régimen previstos en la Sección~\ref{met:metricas} serían necesarios para vincular precisión con bandas de operación.

Los tres horizontes se evalúan como predicciones directas. Esta configuración evita que el error a $1200$ s sea necesariamente la acumulación de sesenta predicciones de $20$ s. Por ello, una curva multihorizonte casi plana tampoco acredita estabilidad recursiva de un modelo de transición. El uso en MPC exige estudiar propagación y sensibilidad a acciones mediante la tarea definida en la Ecuación~\eqref{met:predictor_dinamico}.

La prueba temporal mantiene las mismas $36$ posiciones en entrenamiento y evaluación. En consecuencia, cuantifica pronóstico futuro en ubicaciones conocidas, sin comprobar transferencia a sensores reubicados, puntos interiores o una geometría distinta. La generalización espacial requeriría reservar posiciones o regiones, y conservar la diferencia entre interpolación y extrapolación. Esa distinción afecta la propuesta de reducir instrumentación: un buen error sobre la red completa no demuestra que una subred reproduzca máximos o estados observables con igual suficiencia.
\FloatBarrier

\section{Incertidumbre de la comparación y alcance estadístico}
\label{dis:estadistica}

\subsection{Dependencia temporal y número efectivo de observaciones}

Los $15\,372$ pares posición--origen de la prueba a $1200$ s corresponden a $427$ orígenes compartidos por las $36$ posiciones. La información ambiental y el ciclo determinista son comunes, y los orígenes sucesivos siguen una misma trayectoria. Tratar cada fila como réplica experimental independiente exageraría la precisión inferencial. Los tres días disponibles tampoco representan tres realizaciones meteorológicas independientes: fueron generados mediante las mismas reglas periódicas.

La autocorrelación de orden uno del error cuadrático de persistencia, aproximadamente $0{,}863$ en DS, $0{,}801$ en DHT y $0{,}630$ en humedad, confirma una dependencia relevante para comparar métricas. Esas tres cifras no bastan para calcular un tamaño efectivo de muestra, porque se desconoce la función completa de dependencia y la estructura entre posiciones. La Sección~\ref{met:estadistica} acierta al plantear episodios o bloques temporales como unidades de remuestreo.

La metodología de bloques para observaciones estacionarias conserva dependencia local en lugar de reducir las series a observaciones independientes \cite{kunsch1989}. En este estudio se examinan bloques nominales de $30$, $60$ y $120$ min, mientras que la estacionariedad de los residuos no está establecida. La estabilidad del signo de las mejoras ante esas longitudes aporta sensibilidad favorable dentro de la serie. Para interpretar los intervalos como cobertura inferencial se necesita comprobar que el esquema de remuestreo represente el régimen de residuos y la dependencia entre canales y posiciones.

Con $200$ remuestreos, una cola nominal de $2{,}5\%$ se estima con alrededor de cinco realizaciones en su extremo. Esta resolución limita la estabilidad de los límites percentiles, aunque no invalida la diferencia descriptiva observada. Aumentar remuestreos reduciría esa variabilidad Monte Carlo; disponer de más episodios y regímenes ampliaría la información física. Son mejoras distintas y la primera no sustituye a la segunda.

\subsection{Superioridad frente a persistencia y elección entre ensambles}

Los intervalos positivos de la Tabla~\ref{tab:res_intervalos} respaldan una mejora exploratoria de los ensambles frente a persistencia bajo el procedimiento aplicado. No comparan directamente XGBoost con LightGBM. Restar sus RMSE puntuales o examinar por separado sus intervalos frente a la referencia no produce el intervalo emparejado de la diferencia entre ambos: ese cálculo depende de la covariación de sus errores.

Las diferencias entre ensambles son pequeñas frente a sus errores absolutos. En DS a $1200$ s, la separación es cercana a $0{,}0006\,u_T$ según los valores no redondeados, mientras que los RMSE son aproximadamente $0{,}176\,u_T$. La selección operativa debería basarse en un contraste directo realizado en desarrollo y en criterios previos de complejidad o costo, como establece la Ecuación~\eqref{met:seleccion_ml}. Elegir repetidamente el mínimo de la prueba final convertiría esa prueba en parte de la selección y reduciría su independencia.

La variación interna merece más atención que las diferencias puntuales entre algoritmos. El RMSE DS de XGBoost a $1200$ s alcanza $0{,}3973\,u_T$ en un bloque interno frente a $0{,}1765\,u_T$ en la prueba final (Tabla~\ref{tab:res_validacion_ml}). Esto muestra sensibilidad al tramo de desarrollo y a la cobertura de información disponible. No demuestra por sí solo deriva física ni fuga temporal, pero impide presentar el error final más bajo como prestación uniforme a lo largo de toda la serie.

Finalmente, superioridad predictiva y equivalencia física responden a hipótesis diferentes. Un intervalo positivo de mejora frente a persistencia no verifica que el error esté dentro de una tolerancia de regulación. La equivalencia de la Ecuación~\eqref{met:equivalencia} exige márgenes previos y una estimación de incertidumbre adaptada a la dependencia. Sin unidades térmicas acreditadas, referencias experimentales y márgenes físicos, la aceptación de la planta permanece indeterminada.
\FloatBarrier

\section{PINN y enriquecimiento: coherencia con la física y transferencia pendiente}
\label{dis:hibridacion}

\subsection{Qué significa aproximar la dinámica reducida}

La PINN alcanza errores térmicos entre $5{,}481\times10^{-4}$ y $9{,}741\times10^{-3}$ K respecto de referencias numéricas (Tabla~\ref{tab:res_pinn}). Estas cifras apoyan su capacidad de aproximar el problema reducido ensayado. No pueden compararse como un ranking con los RMSE de sensores, porque cambian fuente, estados, unidades acreditadas, población y tarea. La menor cifra en kelvin no demuestra mayor exactitud del secador que un pronóstico de lecturas expresado en $u_T$.

El uso de ecuaciones como restricción de aprendizaje sigue el planteamiento de las redes informadas por la física \cite{raissi2019}. En esta ejecución, la condición inicial se impone mediante la estructura del mapa condicional de la Ecuación~\eqref{met:flujo_pinn}. El error inicial nulo verifica esa propiedad de construcción. Evaluar el aprendizaje requiere atender al error posterior, a los residuales y a la respuesta bajo estados y entradas reservados.

Un residual pequeño indica cercanía local a las ecuaciones utilizadas. Su relación con el error de trayectoria depende también de la dinámica. Si, en estados normalizados, $\mathbf f$ es Lipschitz con constante $L\geq0$, $\mathbf r=\dot{\widehat{\mathbf x}}-\mathbf f(\widehat{\mathbf x})$ y ambas trayectorias utilizan las mismas entradas, una estimación es
\begin{equation}
\|\mathbf e(t)\|\leq
e^{Lt}\|\mathbf e(0)\|
+\int_0^t e^{L(t-s)}\|\mathbf r(s)\|\,\mathrm ds,
\qquad \mathbf e=\widehat{\mathbf x}-\mathbf x.
\label{eq:dis_error_pinn}
\end{equation}
Esta desigualdad explica por qué residual e inicialización importan, pero no convierte el RMS evaluado en una garantía uniforme. Se necesitarían control del residual sobre el dominio, escalas compatibles y una estimación de amplificación dinámica. Un error de ecuaciones, parámetros o perturbaciones respecto de la planta añadiría contribuciones que esta comparación con el modelo no elimina.

Los pesos uno para datos y física y diez para restricciones no informan por sí solos de la contribución efectiva de cada pérdida. Su interpretación requiere conocer normalización, magnitudes y gradientes. La literatura identifica desequilibrios de gradiente como una dificultad concreta del entrenamiento informado por física \cite{wang2021}. La ausencia de historias de pérdidas y de evaluaciones amplias fuera de los segmentos de ajuste impide atribuir el resultado a un equilibrio óptimo entre términos o a generalización de largo plazo.

La red ODE integrada proporciona además una referencia que debería conservarse al evaluar la utilidad de la PINN. Una sustitución se justificaría por precisión suficiente, latencia, estabilidad multisegmento o conveniencia de diferenciación, con un costo de entrenamiento documentado. Las $24$ referencias por estado no acreditan por sí solas cobertura de estados iniciales, acciones y perturbaciones. El alcance actual es una aproximación numérica condicionada de seis estados, cuya adaptación a los cuatro estados del recinto vacío requiere la correspondencia discutida en la Sección~\ref{dis:alcance}.

\subsection{Generación de datos y efecto del enriquecimiento físico}

La Sección~\ref{res:enriquecimiento} no aporta un contraste base--enriquecido. Por tanto, la mejora de los ensambles sobre la serie sintética no mide una ganancia por añadir CFD o red, y el error de la PINN respecto de una solución numérica tampoco la sustituye. Para atribuir un efecto a la fuente física se necesita conservar algoritmo, información causal y tarea, variando de manera controlada las trayectorias añadidas.

El esquema HPDF de la Sección~\ref{met:hpdf} enfrenta dos riesgos complementarios. Una simulación coherente pero mal parametrizada puede aportar muchas muestras con un mismo sesgo; un conjunto experimental pequeño puede describir adecuadamente la planta en un dominio demasiado estrecho. El número de filas generadas no determina la información independiente. La ponderación por trayectoria de la Ecuación~\eqref{met:perdida_hpdf} ofrece un modo de controlar influencia, cuyo beneficio debe contrastarse sobre observaciones reservadas.

La generación debería concentrarse donde una nueva trayectoria aporte cobertura o reduzca una incertidumbre relevante para decisiones. Repetir ciclos similares bajo los mismos parámetros puede mejorar ajuste al generador sin resolver la discrepancia con la planta. La cavidad, por su geometría cerrada y sus fronteras idealizadas, conserva su función de verificación; incorporarla como si produjera mediciones adicionales del recinto ventilado introduciría una incompatibilidad de tarea.

Una evaluación concluyente del enriquecimiento compararía observaciones solas, observaciones más red y observaciones más proyecciones CFD, con generadores ajustados exclusivamente en desarrollo. La mejora debería examinarse también en conservación, sensibilidad a acciones, error por régimen y costo de generación. El corte disponible permite formular ese contraste y localizar sus requisitos, pero mantiene indeterminado tanto el beneficio como un posible deterioro por transferencia de discrepancias.
\FloatBarrier

\section{Implicaciones para el diseño de control y su implementación}
\label{dis:control}

\subsection{Pronosticar lecturas y evaluar acciones candidatas}

La separación entre pronóstico y transición condicionada por acciones de la Sección~\ref{met:ml} es esencial. Los modelos evaluados predicen canales bajo el régimen de la serie, sin comandos ejecutados y sincronizados. Por ello, sus errores no acreditan el efecto de encender el APK25-5, cambiar el Zeus de bajo a alto o modificar una abertura. Para MPC se necesita una transición cuya respuesta a $\mathbf u$ permita comparar esas decisiones.

En un predictor que no recibe acciones ni un cierre que las relacione con la dinámica, la sensibilidad explícita satisface
\begin{equation}
\frac{\partial\widehat{\mathbf y}_{k+h}}{\partial\mathbf u_{k:k+h-1}}
=\mathbf 0.
\label{eq:dis_sensibilidad_accion}
\end{equation}
En esa formulación, variar una secuencia candidata no altera la salida pronosticada. Un optimizador podría distinguir costos de energía o conmutación, pero no anticiparía su efecto térmico e higrométrico. La ganancia frente a persistencia no resuelve esta carencia de estructura. El predictor podría emplearse en supervisión o pronóstico del régimen observado, mientras se identifica por separado la sensibilidad de la planta.

El ensayo de acciones debe distinguir ganancia térmica, renovación y mezcla. En el recinto vacío, el signo de la respuesta higrométrica depende de $W_{\mathrm{amb}}-W$; el efecto térmico depende también de $T_{\mathrm{amb}}-T_a$, captación y almacenamiento. Una mayor ventilación puede ayudar en un régimen y alejar la referencia en otro. La factibilidad planteada en la Sección~\ref{met:control} debe comprobarse antes de comparar sintonías o arquitecturas.

\subsection{Restricciones de actuación, energía y uniformidad}

El conjunto de la Ecuación~\eqref{met:acciones} contiene seis combinaciones nominales de modos. Ese número no implica seis grados de libertad continuos ni garantiza regulación independiente de temperatura y humedad. Si dos equipos actúan principalmente sobre el mismo mecanismo, la matriz local de respuesta puede tener rango insuficiente o estar mal condicionada. Caudales, aberturas y ensayos de cambio de modo determinarían qué combinaciones ofrecen capacidad adicional.

El control de uniformidad exige otro nivel de información. La red de cuatro estados produce promedios y no identifica por sí sola la posición de un máximo térmico o una zona de baja circulación. Incorporar una penalización espacial sin un cierre contrastado asignaría al controlador una información que su modelo no proporciona. El enlace entre CFD, observación distribuida y red propuesto en la Sección~\ref{met:operadores} es, por tanto, un requisito concreto para regular heterogeneidad.

La energía eléctrica tampoco puede inferirse de las corrientes nominales del Zeus. El producto de tensión y corriente caracteriza potencia aparente en corriente alterna; obtener potencia activa exige el factor de potencia y la condición real de operación. Los $40$ W nominales del APK25-5 permiten una referencia de equipo, pero el consumo por episodio requiere duración efectiva y potencia en el montaje. Sin esas mediciones y sin contenido de agua removida, no se puede calcular consumo específico de secado ni atribuir ahorro a una política.

\subsection{Latencia de predicción y presupuesto del ciclo completo}

Las medianas de unos $5$--$6$ ms para los ensambles son pequeñas respecto del muestreo de $20$ s (Tabla~\ref{tab:res_costo_ml}). En el equipo ensayado, la inferencia individual parece una contribución moderada al presupuesto. El ciclo de control incluye, sin embargo, adquisición, preparación, estimación, escenarios, múltiples evaluaciones, optimización y envío. Una representación del plazo es
\begin{equation}
\begin{split}
\tau_{\mathrm{ciclo}}={}&\tau_{\mathrm{adq/prep}}+\tau_{\mathrm{est}}
+\tau_{\mathrm{esc}}+N_{\mathrm{eval}}\tau_{\mathrm{pred}}\\
&+\tau_{\mathrm{opt,resto}}+\tau_{\mathrm{com}},
\end{split}
\label{eq:dis_presupuesto_control}
\end{equation}
donde $N_{\mathrm{eval}}$ depende del horizonte, escenarios y optimizador, y $\tau_{\mathrm{opt,resto}}$ reúne su costo adicional. La expresión es una organización del presupuesto y no una medición del sistema.

El p99 de $10{,}909$ ms obtenido con diez repeticiones tampoco constituye una garantía del percentil extremo de operación. Con tan pocas observaciones, los cuantiles altos dependen esencialmente de los máximos y de la interpolación utilizada. La ejecución se realizó en una CPU de Windows 11; transferirla a ESP32, un servicio remoto u otro equipo requiere medir latencia, memoria y comunicaciones en esa configuración. El hardware de adquisición no es automáticamente el hardware apropiado para optimización.

La Sección~\ref{res:control} confirma que la comparación en lazo cerrado está aplazada. Por tanto, el resultado actual sustenta la conveniencia de medir el ciclo completo y conserva PID, MRAC y MPC como alternativas formuladas. La selección entre ellas necesita seguimiento, restricciones, energía, conmutaciones y respuesta al desajuste sobre escenarios comunes, conforme a la Sección~\ref{met:comparacion_control}. El error de un predictor constituye una entrada a esa decisión, pero no determina por sí solo el desempeño de la política.
\FloatBarrier

\section{Evaluación de la hipótesis, aportes y prioridades de validación}
\label{dis:objetivos}

\subsection{Respuesta a los objetivos de investigación}

La hipótesis de la Sección~\ref{intro:hipotesis} propone que integrar física, datos, instrumentación y validación permitirá representar con precisión suficiente la dinámica y sustentar el control inteligente. El conjunto actual aporta evidencia favorable sobre la articulación computacional y sobre algunas propiedades de sus componentes. No permite aceptar todavía la suficiencia física y operativa de la cadena completa: faltan una tolerancia definida para la aplicación, observaciones independientes compatibles y resultados de lazo cerrado. Esa insuficiencia tampoco refuta la hipótesis; delimita la parte que continúa abierta.

La Tabla~\ref{tab:dis_objetivos} relaciona los seis objetivos específicos de la Sección~\ref{intro:objetivos} con la evidencia disponible. La valoración distingue el desarrollo de un procedimiento de su comprobación en la planta. De este modo, una metodología bien especificada conserva su aporte sin contabilizar como ejecutadas las etapas que requieren datos o ensayos adicionales.

\begin{table}[htbp]
\centering
\small
\caption{Correspondencia entre objetivos específicos y alcance del corte evaluado.}
\label{tab:dis_objetivos}
\begingroup
\renewcommand{\arraystretch}{1.25}
\setlength{\tabcolsep}{4pt}
\begin{tabularx}{\textwidth}{>{\raggedright\arraybackslash}p{0.30\textwidth}>{\raggedright\arraybackslash}X}
\toprule
Objetivo & Valoración sustentada \\
\midrule
Caracterizar la dinámica y sus variables & Variables y mecanismos definidos; estadísticas de una serie sintética. Caracterización causal y meteorológica de la planta pendiente. \\
Desarrollar modelos físicos y de datos & Implementaciones y pruebas disponibles. El enlace entre el recinto de cuatro estados y las pruebas de seis requiere reconciliación y contraste. \\
Implementar adquisición experimental & Arquitectura, contrato y procesamiento documentados. El archivo generado no acredita disponibilidad, calibración o desempeño de la instalación física. \\
Validar frente a experimentos & Hay verificación numérica y prueba predictiva sintética. La validación independiente de Coffee Solutions sigue abierta. \\
Formular y evaluar control inteligente & Paradigmas, restricciones y protocolo formulados. Su comparación en lazo cerrado no está ejecutada en la exportación. \\
Analizar alcance y transferencia & Se identifican límites de soporte, fuentes, incertidumbre y actuación. La transferencia de parámetros y prestaciones a otros secadores requiere validación propia. \\
\bottomrule
\end{tabularx}
\endgroup
\end{table}

El aporte identificable es una integración que permite someter a examen las interfaces entre componentes: identidad espacial y retardos, magnitudes conservativas y humedad relativa, campos y promedios, predicción y acción. La corrección de antecedentes muestra una mejora verificable de esa cadena. Los balances y los contrastes de integración proporcionan controles cuantitativos. La referencia de persistencia y la separación temporal hacen visible cuándo un algoritmo añade habilidad y cuándo la pierde. Estas contribuciones son más precisas que una afirmación general de validación integral.

\subsection{Prioridades que se desprenden de los resultados}

La primera prioridad es cerrar la correspondencia entre configuración y evidencia. Se necesitan procedencia de unidades y tiempos, geometría contrastada, montaje de sensores y una versión verificable de la red de cuatro estados para el recinto vacío. Esta conciliación precede al ajuste: una calibración numéricamente buena podría compensar una unidad, una posición o una fuente de vapor incorrectas y producir parámetros sin interpretación transferible.

La segunda prioridad es obtener episodios independientes con perturbaciones y acciones sincronizadas. Temperatura y humedad exteriores, presión, irradiancia y modos ejecutados permitirían separar captación, pérdidas y renovación. Las observaciones de superficies apoyarían la estimación de capacidades y conductancias. El diseño debería incluir calentamiento, enfriamiento y cambios de ventilación con duración suficiente para observar respuesta; la repetición de regímenes permitiría cuantificar variabilidad entre episodios.

La tercera prioridad es completar la verificación allí donde las tablas identifican incertidumbre abierta. En la cavidad corresponde examinar oscilación espacial, referencia primaria, detención y balance radiativo total. En el CFD tridimensional hacen falta evolución representativa del episodio, flujos de frontera y refinamientos sobre las magnitudes de contraste. Para la red de cuatro estados corresponde repetir conservación e integración sobre condiciones de uso. Así se evita que la incertidumbre numérica interfiera con la identificación física.

La cuarta prioridad es establecer una jerarquía de predicción sobre una tarea común. Persistencia y una referencia periódica sencilla permitirían evaluar el beneficio adicional de los ensambles; la red identificada proporcionaría el contraste físico. Las variantes híbridas deberían añadir fuentes mediante ablaciones y mantener congelada una prueba experimental externa. La evaluación incluiría cuantiles, regímenes, sensibilidad a acciones y propagación multietapa, además de RMSE global.

La quinta prioridad es comparar control cuando la transición y la capacidad de actuación estén contrastadas. Las políticas recibirían las mismas observaciones, escenarios y restricciones, utilizando una planta de evaluación separada del predictor. El registro del ciclo completo permitiría relacionar seguimiento, consumo, violaciones y respaldo con decisiones efectivamente ejecutadas. Una evaluación posterior con café añadiría gravimetría, temperatura del producto, distribución de carga y calidad, ampliando explícitamente el alcance del recinto vacío.

Estas prioridades también ordenan la discusión sobre costo e instrumentación mínima planteada en la Introducción. Las $36$ posiciones proporcionan un soporte de observación definido, pero no establecen una cantidad óptima de sensores. Reducirlas exige evaluar observabilidad, reconstrucción de extremos, cobertura espacial y tolerancia a fallos sobre datos reservados. Análogamente, seleccionar menos actuadores requiere comprobar que sus modos conservan factibilidad. El ahorro de componentes debe juzgarse junto con la información y la capacidad de regulación que se pierden.

\subsection{Transferibilidad y significado científico del resultado integrado}

Los elementos más transferibles son la estructura de balances, los operadores de observación, la separación temporal de fuentes y el protocolo de comparación. Las capacidades, conductancias, propiedades ópticas, caudales y relaciones espaciales dependen del secador y de su configuración. Los modelos de datos dependen además de la distribución observada. Por ello, adaptar el marco a otro recinto exige reestablecer esas correspondencias y reservar evidencia local para contraste, de acuerdo con la Sección~\ref{subsec:dominio_operacional_validez}.

El estado actual respalda una metodología computacional articulada y parcialmente verificada, con habilidad predictiva demostrada en el caso sintético. También revela que la principal incertidumbre pendiente reside en el vínculo entre modelos, observación física y actuación. Resolver ese vínculo permitirá decidir si la precisión es suficiente para regular el recinto y si la incorporación posterior de producto conserva las prestaciones requeridas. Hasta entonces, las ganancias de predicción, la conservación numérica y la formulación del control constituyen resultados valiosos con alcances propios, cuya delimitación hace comprobable la contribución de la tesis.
\FloatBarrier
